\documentclass[11pt,a4paper]{article}

\usepackage[utf8]{inputenc}
\usepackage[T1]{fontenc}
\usepackage{lmodern}
\usepackage[spanish]{babel}
\usepackage[margin=2.5cm]{geometry}
\usepackage{graphicx}

% Muestra la imagen si existe, si no un recuadro para acordarnos de ponerla
\newcommand{\captura}[2]{%
  \IfFileExists{#1}{\includegraphics[width=\linewidth]{#1}}{\fbox{\parbox{0.9\linewidth}{\centering\vspace{2cm}FALTA: #2\vspace{2cm}}}}%
}

\title{TP D: Desarrollo de aplicaciones distribuidas con BSD sockets}
\author{Grupo 16\\[0.5em] Lautaro Valentín Caminoa, Francisco Krinisky, Maylen Villagran}
\date{Redes y Comunicación, UdeSA\\Noviembre 2026}

\begin{document}
\maketitle

\section{Capturas}

Las capturas se hicieron filtrando por la IP del servidor (\texttt{ip.addr == 167.114.129.206}).

\begin{figure}[h]
  \centering
  \captura{capturas/cp1_wireshark_paquetes.png}{lista de paquetes del checkpoint 1}
  \caption{Checkpoint 1: pedido de autenticación al puerto UDP 20262 y respuesta del servidor.}
\end{figure}

\begin{figure}[h]
  \centering
  \captura{capturas/cp1_wireshark_follow.png}{Follow UDP Stream del checkpoint 1}
  \caption{Checkpoint 1: contenido de los mensajes (Follow UDP Stream). El servidor responde 200 con el puerto TCP 30010.}
\end{figure}

\begin{figure}[h]
  \centering
  % TODO: sacar la captura del desafío 2 y guardarla con este nombre
  \captura{capturas/checkpoint2_wireshark.png}{captura de Wireshark del desafío 2}
  \caption{Desafío 2: autenticación por UDP, conexión TCP al puerto dinámico y reenvío de las PDUs por UDP.}
\end{figure}

\section{Preguntas}

\subsection*{1. ¿Cómo podría mejorar el protocolo para permitir el reenvío de PDUs por UDP con el carácter prohibido 0x7C?}

El \texttt{|} separa los campos del mensaje. Si también aparece dentro de la PDU, el servidor no puede saber dónde termina la PDU. Algunas posibles soluciones podrían ser:

\begin{itemize}
  \item \textbf{Escapar el carácter (byte stuffing):} se elige un carácter de escape, por ejemplo \texttt{\textbackslash}. Dentro de la PDU, \texttt{|} se manda como \texttt{\textbackslash|} y \texttt{\textbackslash} como \texttt{\textbackslash\textbackslash}. El receptor separa los campos solo por los \texttt{|} sin escapar y después saca los escapes. PPP hace lo mismo con su byte de flag (0x7E), usando 0x7D como escape.
  \item \textbf{Codificar la PDU:} mandarla en hexadecimal o en Base64, que no usan el \texttt{|}. El problema es que si bien es más simple, mandar el mensaje crece el doble en hex y un 33\% en Base64.
  \item \textbf{Mandar el largo:} agregar antes de la PDU un campo con su longitud. El receptor lee exactamente esa cantidad de bytes, sin importar qué caracteres traigan.
  \item \textbf{Parsear desde los dos extremos:} los demás campos nunca tienen \texttt{|}: el usuario es alfanumérico, el parse code tiene 3 dígitos y la firma 32 caracteres hex. El servidor podría tomar el usuario hasta el primer \texttt{|}, el parse code y la firma desde el final, y considerar PDU todo lo que queda en el medio. Esto no cambia lo que manda el cliente, solo cómo parsea el servidor.
\end{itemize}

\subsection*{2. ¿Podría el servidor iniciar la conexión TCP contra nuestras computadoras?}

En general no, por el NAT. Nuestras computadoras están en una red privada y el NAT solo agrega una entrada a su tabla cuando un paquete sale desde adentro. La respuesta UDP del servidor nos llega porque contesta un datagrama que mandamos nosotros, y esa entrada ya existe. En cambio, un SYN que llega desde afuera no corresponde a ninguna entrada, así que el NAT no sabe a qué máquina mandarlo y lo descarta. Además, puede haber firewalls que bloqueen las conexiones entrantes, como en la red de la universidad.

Para que funcionara haría falta:
\begin{itemize}
  \item que el cliente haga \texttt{bind}, \texttt{listen} y \texttt{accept} en un puerto y se lo informe al servidor en la autenticación;
  \item tener una IP pública sin NAT, o configurar port forwarding en el router;
  \item que ningún firewall en el camino bloquee la conexión.
\end{itemize}

Es decir, se puede, pero depende de la red de cada uno y en la mayoría de los casos no funcionaría. Si el cliente inicia la conexión, alcanza con que tenga salida a Internet, por eso el protocolo lo hace así.

\subsection*{3. ¿Cómo se vería afectado el funcionamiento de los desafíos ante la pérdida de paquetes?}

En los tramos TCP la pérdida se recupera sola (con demora); en los tramos UDP el mensaje se pierde:

\begin{itemize}
  \item \textbf{Consulta de la IP pública (TCP):} TCP retransmite lo perdido, lo que hace que tarde más.
  \item \textbf{Autenticación (UDP):} no hay retransmisión. Si se pierde el pedido o la respuesta, el cliente espera los 10 segundos y termina sin puerto. Nuestro cliente no reintenta, así que hay que volver a ejecutarlo (o se podría agregar un reintento al vencer el timeout).
  \item \textbf{Conexión al puerto dinámico (TCP):} si se pierde el SYN o el SYN-ACK, TCP lo retransmite. El puerto queda abierto solo 10 segundos, así que si se pierden varios seguidos la conexión puede fallar.
  \item \textbf{Recepción de las PDUs (TCP):} TCP retransmite y entrega los bytes en orden, así que la aplicación no ve la pérdida, solo más demora. Lo que sí puede cambiar es cómo llegan repartidos los bytes en cada \texttt{recv}, y por eso hace falta el framing.
  \item \textbf{Reenvío de las PDUs (UDP):} si se pierde un datagrama, el servidor no recibe esa PDU y el cliente no se entera, porque no hay confirmaciones. También podrían llegar desordenados. Para hacerlo confiable habría que agregar un número de secuencia y un ACK por PDU, y retransmitir si el ACK no llega.
\end{itemize}

\subsection*{4. Dificultades, aprendizajes y consejos}

% TODO: esto es un borrador, cambiarlo con lo que nos pasó de verdad

\textbf{Dificultades}
\begin{itemize}
  \item Entender que en TCP un \texttt{recv} no es un mensaje. Al principio asumíamos que cada \texttt{recv} traía una PDU, y tuvimos que armar el buffer acumulativo.
  \item Implementar MD5 a mano: el padding, el largo en bits al final y que todo es little-endian.
  \item Los strings de C contra los bytes que llegan por la red: \texttt{recv} no agrega el \texttt{\textbackslash 0}, hay que ponerlo.
  \item Los tamaños de los buffers, por ejemplo qué hacer cuando llega una PDU más larga que los 122 bytes.
  \item Coordinar el código entre los tres: separarlo en módulos nos ayudó a trabajar en paralelo.
\end{itemize}

\textbf{Aprendizajes}
\begin{itemize}
  \item La diferencia práctica entre UDP (datagramas, sin garantías, hay que poner timeouts) y TCP (stream confiable, sin límites de mensaje).
  \item Cómo funciona el NAT y por qué hace falta preguntarle a un servicio externo nuestra IP pública.
  \item Que HTTP es solo texto sobre TCP y se puede hablar a mano.
\end{itemize}

\textbf{Consejos para futuros alumnos}
\begin{enumerate}
  \item Familiarizarse con Wireshark desde el principio de la cursada, ya que será necesario para el troubleshooting de este TP.
  \item Si hay algo que no funciona, pueden sniffear el tráfico en wirewshark para encontrar la causa del error.
  \item Nunca asumir que un \texttt{recv} trae un mensaje completo.
  \item Revisar que su implementación del MD5 funcione correctamente con el ejemplo de la consigna antes de mandar algo al servidor.
  \item Verificar que no haya warnings utilizando las flags \texttt{-Wall -Wextra} al compilar.
\end{enumerate}

\end{document}
